\chapter{Literature Review}

This chapter surveys related work along five axes: VRPTW with soft
time windows, safety-aware vehicle routing, multi-objective
$\varepsilon$-constraint methods, quick-commerce and last-mile
routing, and crash-risk calibration in routing. For each axis, we
position the present thesis against the closest prior work and
identify the gap that the SA-VRPTW formulation in Chapter~3
addresses.

\section{VRPTW with Soft Time Windows}

The classical hard-time-window VRPTW \citep{solomon1987} requires
every customer to be served strictly inside a window $[e_i, l_i]$.
The soft variant \citep{koskosidis1992,fu2008} permits late arrival
at a penalty, typically linear in lateness. Convex penalties have
been proposed for SLA-driven settings \citep{taillard1997} but their
linearisation in MILP models is non-trivial and is usually omitted.
Exact column-generation approaches \citep{desrochers1992} solve
small-to-medium instances optimally but do not extend cleanly to the
soft-window-with-convex-penalty regime that q-commerce SLAs demand.

The present thesis adopts $\exp(\beta_{\text{stw}}\tau) - 1$ to
reflect the contractual cost structure of q-commerce, where lateness
becomes operationally expensive faster than linearly (rider
re-routing, customer churn, SLA refunds). The convex term is
piecewise-linearised in the MILP via SOS2 outer envelopes (Chapter~3)
and evaluated exactly in the heuristics, with both code paths sharing
a single objective module.

\section{Safety-Aware Vehicle Routing}

Safety-aware routing originated in hazardous-materials logistics,
where risk was defined as the product of population exposure and
release probability \citep{zografos2004,erkut1998}. More recently,
the principle has been transferred to general urban delivery.
\citet{hoseinzadeh2020} model cumulative crash risk along a tour but
use weighted-sum aggregation, which collapses incompatible quantities
(travel time and expected crash count) and makes the
safety/efficiency trade-off implicit and difficult to interpret.

A second strand penalises hazardous arcs implicitly via inflated
travel times, conflating ``slow'' with ``dangerous'' and giving the
planner no explicit safety dial. The thesis differs in two ways:
first, by adopting an $\varepsilon$-constraint formulation
(Section~\ref{ch3:eps}) so that the per-route crash budget is a
direct knob; second, by separating the network-wide
residential-street budget from the per-rider safety budget, capturing
two operationally distinct concerns (community impact vs.\ rider
survival).

\section{Multi-Objective and \texorpdfstring{$\varepsilon$}{epsilon}-Constraint Methods}

\citet{salmon2023} apply $\varepsilon$-constraint methods to
safety-aware logistics and demonstrate that they dominate weighted-sum
approaches when the planner needs to read off a Pareto frontier
rather than a single weighted optimum. \citet{ehrgott2005} provides
the underlying theoretical framework. The frontier interpretation
matters operationally: a planner picks a target survival probability
$e^{-\bar{R}}$ and the solver returns the cheapest fleet plan that
meets it, or reports infeasibility, with no hidden weight tuning.

This thesis adopts the $\varepsilon$-constraint framing throughout
the empirical study (Chapter~6), generating non-dominated $(F_1,
\bar{R}, \bar{H})$ points by sweeping the budget grid rather than
varying weights.

\section{Quick-Commerce and Last-Mile Routing}

\citet{sinchaisri2023} document the operational intensity of
q-commerce dispatch and the rider-platform tension induced by
piece-rate compensation. Field studies of Indian gig-delivery rider
behaviour \citep{das2020,zografos2004} show how SLA pressure
translates into measurable drops in safety-equipment compliance. The
behavioural Monte Carlo of Chapter~6 is calibrated against these
findings, with a sensitivity sweep documenting how the conclusions
move as the assumed compliance prior shifts.

The Indian q-commerce setting is materially different from the
US/EU last-mile literature in three ways. The street network has
significantly higher residential penetration; the rider population is
exposed two-wheelers rather than cargo bicycles or vans; and the SLA
is a one-shot ten-minute promise rather than a windowed delivery
slot. The customer-placement model used here samples residentially
weighted nodes rather than uniform-disk customers, so the geometric
difficulty of the instances reflects the topology operators face.

\section{Crash-Risk Calibration in Routing}

Most safety-aware routing papers treat the per-edge crash probability
as exogenous, with calibration relegated to a single citation or a
synthetic distribution. This thesis takes a different stance: every
constant in the calibration is traced to a public source, and the
derivation is published as a separate calibration note in the
project repository. \citet{morth2022} is the headline aggregate
source; \citet{mohan2017} provides the directional ordering
(arterials more dangerous than collectors more dangerous than local
streets); and the bounded local proxy uses graph structure available
through OSMnx \citep{boeing2017osmnx}.

The deliberate constraint is that the calibration chain does not
require access to the protected Integrated Road Accident Database
(iRAD) microdata. The price is that the per-edge probability is a
calibrated prior rather than a measurement; the benefit is full
public-source reproducibility.

\section{Heuristics for Capacitated Multi-Depot Routing}

Two metaheuristic strands inform the implementation of Chapter~5.
Genetic algorithms for VRP \citep{holland1975} routinely use the
giant-tour-with-split decoder \citep{prins2004}, in which a single
permutation chromosome over customers is decoded into routes by
solving a shortest-path problem on an auxiliary acyclic graph. This
keeps the search space free of capacity-/time-feasibility scaffolding
and pushes feasibility into the deterministic decoder.

The Adaptive Large Neighbourhood Search of
\citet{ropke2006} alternates destroy and repair operators selected
under an adaptive roulette-wheel scheme, with simulated-annealing
acceptance for diversification. ALNS is particularly attractive for
SA-VRPTW because tightening an $\varepsilon$-budget often requires
reshaping several short routes simultaneously, which is exactly the
move large-neighbourhood operators are designed to make.

\section{Research Gap}

The literature contains separate strands on (i) VRPTW with soft time
windows, (ii) safety-aware routing with crash-risk objectives, (iii)
$\varepsilon$-constraint multi-objective optimisation, (iv)
quick-commerce dispatch, and (v) public-source road-safety
calibration. No single existing study combines all five into a
unified, reproducible framework grounded in real Indian dark-store
coordinates and a documented calibration chain. The contribution of
this thesis is to do exactly that: a calibration-honest SA-VRPTW
formulation with three solvers under one evaluation contract,
released for full external audit.
